\chapter{Apraxia}

\section*{Definition}
Apraxia is a disorder of skilled, purposive movement that cannot be attributed to weakness, sensory loss, ataxia, incoordination, language comprehension deficits, or you lack of cooperation. People demonstrate difficulty performing learned motor tasks on command despite intact motor and sensory systems.

\section*{What Happens in Your Body}
Apraxia results from dysfunction of the neural networks that store and retrieve motor engrams -- the learned spatial and temporal representations of complex movements. The dominant parietal lobe (especially the supramarginal gyrus) integrates sensory feedback with motor plans, while the supplementary motor area and premotor cortex sequence these plans into action. Lesions disrupting these networks disconnect the conceptual representation of a movement from its execution, or the execution plan from the primary motor cortex.

\section*{Causes}
\begin{itemize}
 \item Stroke: Left parietal lobe infarction (especially superior parietal lobule and supramarginal gyrus), anterior cerebral artery infarction affecting the supplementary motor area, subcortical lesions of the basal ganglia or thalamus
 \item Neurodegenerative disease: Corticobasal degeneration (asymmetric apraxia is a hallmark), Alzheimers disease (ideational apraxia emerges in moderate stages), progressive supranuclear palsy, dementia with Lewy bodies
 \item Traumatic brain injury: Contusions or diffuse axonal injury involving frontal and parietal association cortices
 \item Neoplastic: Gliomas or metastases involving the dominant parietal lobe or premotor frontal cortex
 \item Infectious: Brain abscess, herpes simplex encephalitis involving frontoparietal regions
 \item Toxic-metabolic: Carbon monoxide poisoning, hypoxic-ischemic injury, Wilson disease (basal ganglia involvement)
 \item Demyelinating: Multiple sclerosis plaques affecting white matter tracts between parietal and frontal regions
\end{itemize}

\section*{When to See a Doctor}

\begin{itemize}
 \item Sudden-onset difficulty performing familiar tasks (using utensils, dressing, waving) in the absence of weakness
 \item Apraxia accompanied by aphasia, hemispatial neglect, or hemiparesis (suggests dominant hemisphere stroke)
 \item Progressive decline in ability to sequence multi-step tasks with preserved strength and sensation
 \item Any apraxia with new headache, seizure, or altered mental status
\end{itemize}

\section*{Related Symptoms}
\begin{itemize}
 \item Aphasia (especially Broca or conduction aphasia)
 \item Agraphia
 \item Alexia
 \item Hemispatial neglect
 \item Ideomotor apraxia (impaired pantomime of tool use)
 \item Ideational apraxia (impaired multi-step task sequencing)
 \item Limb-kinetic apraxia (clumsy, coarse distal movements)
 \item Orofacial apraxia (inability to perform facial movements on command)
 \item Dysarthria
 \item Alien limb phenomenon
\end{itemize}
\textbf{Important:} Limb apraxia is often overlooked during bedside examination because compensatory strategies mask the deficit; testing should include transitive gestures (e.g., hammering, brushing teeth) elicited by verbal command, imitation, and with the actual tool.

\section*{Self-Care / Home Management}
\begin{itemize}
 \item Maintain a structured daily routine with familiar objects placed consistently to reduce reliance on motor planning.
 \item Use verbal self-cueing strategies (narrating steps aloud) during multi-step tasks such as dressing or meal preparation.
 \item Remove hazards such as sharp knives or hot surfaces if safety is compromised by task sequencing errors.
 \item Encourage you to slow down and focus on one step at a time; occupational therapy referral is essential for adaptive strategy training.
\end{itemize}

\section*{How Your Doctor Will Evaluate You}
\begin{itemize}
 \item Comprehensive neurological examination including strength, sensation, coordination, and language assessment to rule out alternative explanations.
 \item Apraxia-specific testing: transitive and intransitive gesture pantomime, imitation of meaningless hand postures, tool use with real objects, multi-step task performance.
 \item Neuroimaging (MRI brain with and without contrast) to identify structural lesions, atrophy patterns, or white matter disease.
 \item Neuropsychological evaluation for suspected neurodegenerative cause when stroke is excluded and symptoms progress.
\end{itemize}

\section*{Treatment Options}
\begin{itemize}
 \item Stroke rehabilitation with occupational and speech therapy focusing on gesture training, constraint-induced therapy, and chaining of multi-step tasks.
 \item Treat underlying cause: thrombolysis or thrombectomy for acute ischemic stroke, immunosuppression for autoimmune encephalitis, surgical resection for tumor.
 \item Pharmacotherapy for associated conditions (e.g., donepezil for Alzheimers, levodopa for corticobasal syndrome) may slow progression but does not directly treat apraxia.
 \item Environmental modification and caregiver education regarding cueing strategies, simplification of tasks, and realistic goal-setting for progressive cases.
\end{itemize}

\section*{References}
\begin{thebibliography}{99}
\bibitem{apraxia:ref1} Wheaton LA, Hallett M. ``Ideomotor apraxia: a review.'' \textit{J Neurol Sci}. 2007;260(1-2):1-10.
\bibitem{apraxia:ref2} Gross RG, Grossman M. ``The neuropsychology of apraxia.'' \textit{J Clin Exp Neuropsychol}. 2008;30(7):769-781.
\bibitem{apraxia:ref3} Heilman KM, Watson RT, et al. ``Apraxia.'' In: \textit{Bradley\'s Neurology in Clinical Practice}, 8th ed. Elsevier, 2021.
\bibitem{apraxia:ref4} Miller BL, Boeve BF, eds. \textit{The Behavioral Neurology of Dementia}, 2nd ed. Cambridge University Press, 2016.
\bibitem{apraxia:ref5} UpToDate. ``Apraxia: clinical features and evaluation.'' Wolters Kluwer, 2023.
\end{thebibliography}
